\documentclass[12pt]{article}
\usepackage[utf8]{inputenc}
\usepackage{bbm}
\usepackage{geometry}
\geometry{a4paper,scale=0.75}
\linespread{1.5}
\usepackage{graphicx} 
\usepackage{float} 
\usepackage{subcaption} 
\usepackage{enumerate}
\usepackage{amsmath}
\usepackage{array}
\newcolumntype{P}[1]{>{\raggedright\arraybackslash}p{#1}}
\usepackage{booktabs}
\usepackage{multirow}
\usepackage{amsfonts}
\usepackage[english]{babel}
\usepackage{amsthm}
\usepackage{dcolumn}
\usepackage{multicol}
\usepackage{stfloats}
\usepackage{placeins}
\usepackage{lscape}
\usepackage[figuresright]{rotating}
\RequirePackage{pdflscape}
\usepackage[toc,page]{appendix}
\usepackage{geometry}
\usepackage{longtable}
\usepackage{comment}
\usepackage{xcolor}
% \usepackage{amsmath}
\usepackage{amssymb}
\usepackage{empheq}
\usepackage{newtxtext,newtxmath}

% -------- enumerated sub-labels (a), (b), … --
\usepackage{enumitem}
\setlist[enumerate,1]{label=(\alph*),ref=\alph*}
% ---------------------------------------------

\usepackage{hyperref}
\hypersetup{hidelinks,
	colorlinks=true,
	allcolors=black,
	pdfstartview=Fit,
	breaklinks=true}
\usepackage{csquotes}
\usepackage{natbib}
\newtheorem{definition}{Definition}
\newtheorem{theorem}{Theorem}
%\newtheorem{proposition}[theorem]{Proposition}
\newtheorem{proposition}{Proposition}[section]
%\newtheorem{lemma}[theorem]{Lemma}
\newtheorem{lemma}{Lemma}[section]
%\newtheorem{corollary}[theorem]{Corollary}
\newtheorem{corollary}{Corollary}[section]
\newtheorem*{remark}{Remark}
\newtheorem{example}{Example}
\newtheorem{exercise}{Exercise}[section]
\numberwithin{equation}{section}
\newtheorem{assumption}{Assumption}[section] % number within sections

\title{Production Networks and Formation of Aggregate Uncertainty}
\author{
Harlly Zhou\footnotemark[1]\\\texttt{hzhouci@ust.hk}
}
\date{\today}

\begin{document}

\begin{titlepage}
    \begin{center}
        {\LARGE\bfseries Uncertainty, Production Networks, and Optimal Monetary Policy\par}
        \vspace{1.5em}
        {\large
        Harlly Zhou\footnotemark[1] (\texttt{hzhouci@ust.hk})\\
        }
        \vspace{1em}
        \today
    \end{center}
    \vspace{2em}
    \begin{abstract}
        \sloppy
        \noindent 
         \vspace{1em}

        \noindent\textbf{JEL Codes:} 


        \noindent\textbf{Keywords:} 
        \fussy
    \end{abstract}
    \thispagestyle{empty}
    \setcounter{page}{0}
\end{titlepage}


\section{Benchmark Model}\label{sec:benchmark_model}
\subsection{Model Setup}
\subsubsection{Sectors and Firms}

Consider a closed-economy with $N$ industries, indexed by $i=1,\cdots,N$. Each industry has a representative firm that produces a differentiated good $y_{i,t}$ using capital $k_{i,t}$, labor $n_{i,t}$, and intermediate inputs from industry $j$, $\{x_{ij,t}\}_{j=1}^{N}$. The production function is given by
\begin{align}\label{eq:prod_func}
    y_{i,t}
    =
    z_{i,t}\zeta_i
    \left(
        k_{i,t}^{\eta_i}n_{i,t}^{1-\eta_i}
    \right)^{a_i}
    \left[
        \prod_{j=1}^{N}
        x_{ij,t}^{\alpha_{ij}}
    \right]^{1-a_i},
\end{align}
where $z_{i,t}$ is the productivity of industry $i$, $\zeta_i$ is a normalizing factor, and $x_{ij,t}$ is the quantity of good $j$ used by industry $i$ as an intermediate input. Assume that $\sum_{j=1}^{N}\alpha_{ij}=1$.

Capital evolves according to
\begin{align}\label{eq:capital_accumulation}
    k_{i,t+1}
    =
    \xi_{i,t}
    \left[
        (1-\delta)k_{i,t}+i_{i,t}
    \right],
\end{align}
where $\delta$ is the physical capital depreciation rate, $i_{i,t}$ is investment, and $\xi_{i,t}\in[0,1]$ is the remaining proportion of renewed capital after potential liquidation. For the benchmark, we assume that there is no adjustment cost for investment.

The stochastic productivity process is an $AR(1)$ process:
\begin{align}\label{eq:productivity_process}
    \log z_{i,t}
    = \rho^z_{i} \log z_{i,t-1}
    +\sigma_{i,t}\epsilon^z_{i,t},
\end{align}
where $\rho^z_{i}\in(0,1)$ is the persistence of productivity. $\sigma_{i,t}^2$ is the innovation variance, and $\epsilon^z_{i,t}$ is the productivity innovation. Denote $\boldsymbol \sigma_t = (\sigma_{1,t},\ldots,\sigma_{N,t})'\in \mathbb{R}^N$. Suppose that the innovation variance of each sector moves over time according to an $AR(1)$ process: 
\begin{align}\label{eq:innovation_variance_process}
    \log \sigma_{i,t} = \rho^\sigma_{i} \log \sigma_{i,t-1} + \epsilon^\sigma_{i,t},
\end{align}
where $\rho^\sigma_{i}\in(0,1)$ is the persistence of the innovation variance. Let
\begin{align*}
    \boldsymbol\epsilon^z_t
    &=
    (\epsilon^z_{1,t},\ldots,\epsilon^z_{N,t})'\in \mathbb{R}^N,
    \qquad
    \boldsymbol\epsilon^z_t
    \sim
    \mathcal N(\mathbf 0, \mathbf{I_N}), \\
    \boldsymbol\epsilon^\sigma_t
    &=
    (\epsilon^\sigma_{1,t},\ldots,\epsilon^\sigma_{N,t})'\in \mathbb{R}^N,
    \qquad
    \boldsymbol\epsilon^\sigma_t
    \sim
    \mathcal N(\mathbf 0, \mathbf{\Sigma^\sigma}),
\end{align*}
where $\mathbf{I_N}$ is the $N\times N$ identity matrix and $\mathbf{\Sigma^\sigma}$ is the $N\times N$ variance-covariance matrix. For each $i$, $\epsilon^z_{i,t}$ and $\epsilon^\sigma_{i,t}$ have mean zero and unit variance. $\epsilon^\sigma_{i,t}$ realizes together with the innovation variance $\sigma_{i,t}$ and $\epsilon^z_{i,t}$ realizes together with the productivity $z_{i,t}$.



At the beginning of each period $t$, the innovation variance realizes, while the productivity innovation and thus productivity have not yet realized. The firm chooses borrowing $b_{i,t}$ from households to maximize its lifetiem value. Borrowing is used to finance labor costs, investment costs, and the costs of intermediate inputs. Investment $i_{i,t}$ is simultaneously chosen at this stage.

After productivity realizes, prices are flexible to clear the goods market. Given the borrowing and investment amount chosen in the first stage, firms choose labor demand $n_{i,t}$ and intermediate inputs $x_{i1,t},\ldots,x_{iN,t}$.

Goods-market clearing in industry $i$ is
\begin{align}\label{eq:goods_market_clearing}
    y_{i,t}
    =
    c_{i,t}+\sum_{j=1}^{N}x_{ji,t}.
\end{align}
Therefore, after sales are realized, the firm earns nominal revenue $p_{i,t} y_{i,t}$, subject to a proportional revenue tax $\tau_{i,t}$. Here we assume that revenue is observable and that creditors have seniority over revenue and pledgeable assets.

The firm repays borrowing $b_{i,t}$ at the interest rate $r_{i,t}$, where $r_{i,t}$ is the sector-specific interest rate set by creditors that is higher than the risk-free rate to compensate for the default risk. Default occurs if the firm does not have sufficient enforceable resources. If after-tax sales revenue is insufficient for repayment, the firm liquidates some of the renewed capital
\[
    \tilde{k}_{i,t} \equiv (1-\delta)k_{i,t}+i_{i,t}.
\]
For simplicity, assume that there is a maximum proportion of capital liquidation, so that at least $\underline{\xi}\in(0,1)$ of renewed capital is carried into the next period. Therefore, $\xi_{i,t}\in[\underline{\xi},1]$. This restriction abstracts from firm entry and bankruptcy. If there is a positive remaining profit, the firm distributes it to households.

In summary, firms make decisions in two stages. After the innovation variances are realized but before productivity is realized, firm $i$ chooses borrowing $b_{i,t}$ and investment $i_{i,t}$. Its first-stage problem is
\begin{align}\label{eq:firm_first_stage_decision}
    J
    \left(k_{i,t}; \boldsymbol\sigma_t\right)
    =
    \max_{b_{i,t}, i_{i,t}}
    \mathbb{E}_{i,t}
    \Bigg[
            d_{i,t}
            +
            \mathcal{M}_{t,t+1}
            J\left(k_{i,t+1}; \boldsymbol\sigma_{t+1}\right)
    \Bigg],
\end{align}
subject to the capital LoM \eqref{eq:capital_accumulation} and the working capital constraint $b_{i,t} \geq w_t n_{i,t} + \sum_{j=1}^{N} p_{j,t} x_{ij,t} + p_{K,t} i_{i,t}$. Here $\mathbb{E}_{i,t}$ denotes the expectation conditional on the information available after the innovation variances are realized but before productivity is realized. $d_{i,t}$ denotes the state-contingent maximized equity payout of the firm after productivity is realized, defined by
\begin{align*}
    d_{i,t}
    =
    \max\left\{
        0,
        (1-\tau_{i,t})p_{i,t} y_{i,t}
        -
        (1+r_{i,t})b_{i,t}
    \right\},
\end{align*}
whose realized value is determined in a sate-contingent manner in the second stage.

In the second stage, the firm chooses inputs after productivity is realized:
\begin{align}\label{eq:firm_second_stage_decision}
    \min_{n_{i,t},\{x_{ij,t}\}_{j=1}^{N}}
    w_t n_{i,t}
    + \sum_{j=1}^{N} p_{j,t} x_{ij,t}
\end{align}
subject to
\begin{align*}
    b_{i,t} &\geq w_t n_{i,t}
    + \sum_{j=1}^{N} p_{j,t} x_{ij,t}
    + p_{K,t} i_{i,t},
    \\
    y_{i,t}
    &=
    z_{i,t}\zeta_i
    \left(
        k_{i,t}^{\eta_i}n_{i,t}^{1-\eta_i}
    \right)^{a_i}
    \left[
        \prod_{j=1}^{N}
        x_{ij,t}^{\alpha_{ij}}
    \right]^{1-a_i},
\end{align*}
where $w_t$ is the nominal wage, $p_{j,t}$ is the price of good $j$, $p_{K,t}$ is the price of physical capital. We can get the following expression in log form for the marginal cost:
\begin{align}\label{eq:marginal_cost}
    \text{mc}_{i,t} = \frac{1}{(1-\eta_i)a_i} \left[(1-\eta_i)a_i \text{w}_t + (1-a_i)\sum_{j=1}^N \alpha_{ij} \text{p}_{j,t} - \eta_i a_i \text{k}_{i,t} - \text{z}_{i,t}\right],
\end{align}
where the normalization factor is set as
\begin{align*}
    \zeta_i = \dfrac{[1-\eta_i a_i]^{1-\eta_i a_i}}{[(1-\eta_i)a_i]^{1-\eta_i a_i}\left[\prod_{j=1}^N \alpha_{ij}^{\alpha_{ij}}\right]^{1-a_i}}.
\end{align*}

For the first-stage problem, the firm's FOCs are
\begin{align}\label{eq:first_stage_focs}
    [b_{i,t}]: \quad &\mathbb{E}_{i,t}\left[\lambda^{f}_{i,t} - (1+r_{i,t})\mathbf{1}_{\Xi^1_{i,t}} - \mathcal{M}_{t,t+1} \frac{1+r_{i,t}}{p_{K,t}} J'(k_{i,t+1})\mathbf{1}_{\Xi^2_{i,t}} \right] = 0, \\
    [i_{i,t}]: \quad &\mathbb{E}_{i,t}\left[\lambda^{f}_{i,t}p_{K,t} - \mathcal{M}_{t,t+1} J'(k_{i,t+1})\left(\mathbf{1}_{\Xi^1_{i,t} \cup \Xi^2_{i,t}} + \underline{\xi}\mathbf{1}_{\Xi^3_{i,t}}\right)\right] = 0,
\end{align}
where $\mathbf{1}$ denotes the indicator function of the sets $\Xi^1_{i,t}, \Xi^2_{i,t}, \Xi^3_{i,t}$. The sets are defined based on the three possible states of the productivity innovation:
\begin{align*}
    \Xi^1_{i,t} &:= \left\{z_{i,t}: (1-\tau_{i,t})p_{i,t} y_{i,t} \geq (1+r_{i,t})b_{i,t}\right\},\\
    \Xi^2_{i,t} &:= \left\{z_{i,t}: 0 < (1+r_{i,t})b_{i,t} - (1-\tau_{i,t})p_{i,t} y_{i,t} < (1-\underline{\xi}) p_{K,t} \tilde{k}_{i,t}\right\},\\
    \Xi^3_{i,t} &:= \left\{z_{i,t}: (1+r_{i,t})b_{i,t} - (1-\tau_{i,t})p_{i,t} y_{i,t} \geq (1-\underline{\xi}) p_{K,t} \tilde{k}_{i,t}\right\}.
\end{align*}
The Envelope condition is
\begin{align*}
    J'(k_t) = &\mathbb{E}_{i,t}\bigg[\eta_i a_i(1-\tau_{i,t})p_{i,t}\frac{y_{i,t}}{k_{i,t}}\mathbf{1}_{\Xi^1_{i,t}} \\
    &+ (1-\delta)\mathcal{M}_{t,t+1} J'(k_{t+1})\left(\mathbf{1}_{\Xi^1_{i,t} \cup \Xi^2_{i,t}} + \underline{\xi}\mathbf{1}_{\Xi^3_{i,t}} + \eta_i a_i(1-\tau_{i,t})p_{i,t}\frac{y_{i,t}/k_{i,t}}{p_{K,t}}\mathbf{1}_{\Xi^2_{i,t}}\right)\bigg].
\end{align*}
Equivalently, let $f_{i,t}(z)$ denote the conditional density of $z_{i,t}$ given the information available at the first stage. Since $\log z_{i,t}=\rho_i^z\log z_{i,t-1}+\sigma_{i,t}\epsilon^z_{i,t}$ and $\epsilon^z_{i,t}\sim\mathcal N(0,1)$,
\begin{align*}
    f_{i,t}(z)
    =
    \frac{1}{z\sigma_{i,t}}
    \phi\left(
        \frac{\log z-\rho_i^z\log z_{i,t-1}}{\sigma_{i,t}}
    \right),
\end{align*}
where $\phi(\cdot)$ is the standard normal density. Define the conditional probabilities
\begin{align*}
    \pi^m_{i,t}
    :=
    \Pr_{i,t}\left(z_{i,t}\in\Xi^m_{i,t}\right)
    =
    \int_{\Xi^m_{i,t}} f_{i,t}(z)\,dz,
    \qquad
    m=1,2,3,
\end{align*}
and $\pi^{12}_{i,t}:=\pi^1_{i,t}+\pi^2_{i,t}$. Then the first-order conditions can be written explicitly as
\begin{align}\label{eq:first_stage_focs_probability}
    [b_{i,t}]: \quad
    &\mathbb{E}_{i,t}\left[\lambda^f_{i,t}\right]
    -
    (1+r_{i,t})\pi^1_{i,t}
    -
    \frac{1+r_{i,t}}{p_{K,t}}
    \pi^2_{i,t}
    \mathbb{E}_{i,t}\left[
        \mathcal M_{t,t+1}
        J'(k_{i,t+1};\boldsymbol\sigma_{t+1})
        \,\middle|\,
        z_{i,t}\in\Xi^2_{i,t}
    \right]
    =
    0,
    \\
    [i_{i,t}]: \quad
    &p_{K,t}\mathbb{E}_{i,t}\left[\lambda^f_{i,t}\right]
    -
    \pi^{12}_{i,t}
    \mathbb{E}_{i,t}\left[
        \mathcal M_{t,t+1}
        J'(k_{i,t+1};\boldsymbol\sigma_{t+1})
        \,\middle|\,
        z_{i,t}\in\Xi^1_{i,t}\cup\Xi^2_{i,t}
    \right]
    \notag\\
    &\qquad
    -
    \underline{\xi}\pi^3_{i,t}
    \mathbb{E}_{i,t}\left[
        \mathcal M_{t,t+1}
        J'(k_{i,t+1};\boldsymbol\sigma_{t+1})
        \,\middle|\,
        z_{i,t}\in\Xi^3_{i,t}
    \right]
    =
    0.
\end{align}

\subsubsection{Households}

There are representative households with unit mass supplying labor and credit to firms and choosing consumption quantities $\{c_{i,t}\}_{i=1}^{N}$. Household utility is given by
\begin{align}\label{eq:hh_utility}
    U(C_t,N_t)
    =
    \log C_t
    -
    \chi\frac{N_t^{1+\varphi}}{1+\varphi},
\end{align}
where consumption is an aggregator of consumption goods,
\begin{align*}
    C_t
    :=
    \left[
        \sum_{i=1}^{N}
        c_{i,t}^{
            \frac{\theta-1}{\theta}
        }
    \right]^{
        \frac{\theta}{\theta-1}
    },
\end{align*}
where $\theta\in(0,+\infty)$ is the elasticity of substitution between consumption goods, and $N_t=\sum_{i=1}^{N}n_{i,t}$ is total labor supply.

The household budget constraint is
\begin{align}\label{eq:hh_budget_constraint}
    &w_tN_t
    +\sum_{i=1}^{N}
        (1+\tilde{r}_{i,t})b_{i,t}
    +\sum_{i=1}^N e_{i,t} (d_{i,t} + q_{i,t}) + T_t \notag\\
    &\qquad \qquad\geq
    \sum_{i=1}^{N}
        p_{i,t}c_{i,t}
    +\sum_{i=1}^{N}b_{i,t+1} +\sum_{i=1}^N e_{i,t+1}  q_{i,t}.
\end{align}
Here $e_{i,t}$ is the number of shares of firm $i$'s equity holdings, $q_{i,t}$ is the price of firm $i$'s equity, and $\tilde{r}_{i,t}$ is the realized return on borrowing. It is defined by
\begin{align}\label{eq:hh_realized_return}
    (1+\tilde{r}_{i,t})b_{i,t} \equiv
    \min\Bigg\{
        (1+r_{i,t})b_{i,t},
        (1-\tau_{i,t})p_{i,t} y_{i,t}
        +
        (1-\underline{\xi})p_{K,t}
        \left[
            (1-\delta)k_{i,t}+i_{i,t}
        \right]
    \Bigg\},
\end{align}
where $D_t:= \sum_{i=1}^N d_{i,t}$ is total nominal profit from firms, and $T_t$ is a lump-sum government transfer. Total consumption expenditure is
\begin{align*}
    P_tC_t
    :=
    \sum_{i=1}^{N}
        p_{i,t}c_{i,t},
\end{align*}
where $P_t$ is the aggregate price index.

In summary, the household problem is
\begin{align}\label{eq:hh_problem}
    \max_{\left\{c_{i,t}, n_{i,t}, e_{i,t}, r_{i,t}\right\}_{i=1}^{N}} \mathbb{E}_t \sum_{t=0}^{+\infty} \beta^t U(C_t, N_t)
\end{align}
subject to the budget constraint \eqref{eq:hh_budget_constraint}.

\subsubsection{Government}
There is a government that conducts fiscal and monetary policy. On the fiscal side, the government collects taxes or provides subsidies to firms and rebates the resulting revenue to households as a lump-sum transfer:
\begin{align*}
    T_t
    \equiv
    \sum_{i=1}^{N}
        \tau_{i,t}p_{i,t}y_{i,t}.
\end{align*}








\clearpage
\bibliographystyle{aer}
\bibliography{network_uncertainty}


\clearpage


%%%%%%%%%%%%%%%%%%%%%%%%%%
% Publication Appendix
%%%%%%%%%%%%%%%%%%%%%%%%%%
%%%%%%%%%%%%%%%%%%%%%%%%%%%%%%%%%%%%%%%%%%%%%%%%%%%%%%%%%%%%%%%%%%%%%%%%%%%%%%%%%%%%%%%
%%%%%%%%%%%%%%%%%%%%%%%%%%%%%%%%%%%%%%%%%%%%%%%%%%%%%%%%%%%%%%%%%%%%%%%%%%%%%%%%%%%%%%%
%%%%%%%%%%%%%%%%%%%%%%%%%%%%%%%%%%%%%%%%%%%%%%%%%%%%%%%%%%%%%%%%%%%%%%%%%%%%%%%%%%%%%%%
%%%%%%%%%%%%%%%%%%%%%%%%%%%%%%%%%%%%%%%%%%%%%%%%%%%%%%%%%%%%%%%%%%%%%%%%%%%%%%%%%%%%%%%
% \clearpage
\begin{appendix}
% Formatting options
\addcontentsline{toc}{section}{Appendices}
\renewcommand{\thesection}{\Roman{section}} 
\renewcommand{\thesubsection}{\Roman{subsection}}
\renewcommand{\theequation}{\thesection\arabic{equation}}
\numberwithin{equation}{section}
% \setcounter{page}{1}
\setcounter{figure}{0} \renewcommand{\thefigure}{\Roman{section}.\arabic{figure}}
\setcounter{table}{0} \renewcommand{\thetable}{\Roman{section}.\arabic{table}}
\setcounter{subsection}{0} \renewcommand{\thesubsection}{\Roman{section}.\arabic{subsection}}
\setcounter{equation}{0} \renewcommand{\theequation}{\Roman{section}.\arabic{equation}}
\setcounter{theorem}{0} \renewcommand{\thetheorem}{\Roman{section}.\arabic{theorem}}
\setcounter{definition}{0} \renewcommand{\thedefinition}{\Roman{section}.\arabic{definition}}
% \setcounter{proposition}{0} \renewcommand{\theproposition}{\Roman{section}.\arabic{proposition}}
\setcounter{assumption}{0} \renewcommand{\theassumption}{\Roman{section}.\arabic{assumption}}
% \setcounter{hyp}{0} \renewcommand{\thehyp}{\Roman{section}.\arabic{hyp}}
\setcounter{footnote}{0}
%%%%%%%%%%%%%%%%%%%%%%%%%%%%%%%%%%%%%%%%%%%%%%%%%%%%%%%%%%%%%%%%%%%%%%%%%%%%%%%%%%%%%%%
%%%%%%%%%%%%%%%%%%%%%%%%%%%%%%%%%%%%%%%%%%%%%%%%%%%%%%%%%%%%%%%%%%%%%%%%%%%%%%%%%%%%%%%
%%%%%%%%%%%%%%%%%%%%%%%%%%%%%%%%%%%%%%%%%%%%%%%%%%%%%%%%%%%%%%%%%%%%%%%%%%%%%%%%%%%%%%%
%%%%%%%%%%%%%%%%%%%%%%%%%%%%%%%%%%%%%%%%%%%%%%%%%%%%%%%%%%%%%%%%%%%%%%%%%%%%%%%%%%%%%%%



\end{appendix}

\end{document}